% !TEX program = xelatex
% ============================================================
% SLB 技术文档：6.5 物理层通用功能（整章整合）
% 规范：slb-doc-generator / tech-doc-style-guide.md
% 模块类型：通用功能（章级索引 + 端到端）
% 分册：6.5.1-6.5.3 / 6.5.4-6.5.5 / 6.5.6 / 6.5.7-6.5.9
% ============================================================
\documentclass[11pt,a4paper]{ctexart}

\usepackage{amsmath,amssymb,amsfonts}
\usepackage{booktabs}
\usepackage{geometry}
\usepackage{hyperref}
\usepackage{enumitem}
\usepackage{xcolor}
\usepackage{longtable}
\usepackage{array}

\geometry{left=2.5cm,right=2.5cm,top=2.5cm,bottom=2.5cm}
\hypersetup{colorlinks=true,linkcolor=blue,urlcolor=blue}

\newcolumntype{L}[1]{>{\raggedright\arraybackslash}p{#1}}

\title{%
  \textbf{星闪 SLB 物理层}\\[0.4em]
  \Large 6.5 物理层通用功能技术文档（整合）\\[0.3em]
  \large CRC/序列 \textbullet\ 功控测量 \textbullet\ 调制 \textbullet\ 加扰/SC-FDMA/ZC
}
\author{依据 T/XS 10001-2025《星闪无线通信系统 接入层\\
同步低时延宽带空口 SLB 技术要求和测试方法》整理\\
（团体标准 20250326 版）}
\date{2026 年 7 月 27 日}

\begin{document}
\maketitle
\tableofcontents
\newpage

%======================================================================
\part{总览}
%======================================================================

\section{文档范围与模块划分}

本文档为 T/XS 10001-2025 第 6.5 章``物理层通用功能 / 共性基础算法''的\textbf{整章整合入口}：给出 6.5.1\textasciitilde 6.5.9 模块矩阵、跨节端到端流水线、统一约束与分册索引。公式级细节以四份分册为准，避免多处复制导致漂移。

\begin{longtable}{L{2.4cm}L{3.5cm}L{4.5cm}L{3.5cm}}
\caption{6.5 章模块划分与分册\label{tab:65-scope}} \\
\toprule
\textbf{节号} & \textbf{主题} & \textbf{核心行为} & \textbf{分册文件} \\
\midrule
6.5.1 & CRC & 五种 $g(D)$ 编码/校验 & \texttt{6.5.1-6.5.3.tex} \\
6.5.2 & Gold 序列 & 31 阶双 $m$ + 1600 偏移 & 同上 \\
6.5.3 & 伪随机 QPSK & 161 点 RS；$k{=}80$ 置零 & 同上 \\
6.5.4 & 功率控制 & $P{=}10\log_{10}M{+}P_o{+}P_L$；限幅 & \texttt{6.5.4-6.5.5.tex} \\
6.5.5 & 测量量 & RSRP/RSSI/RSRQ/SINR/CBRA/PMS & 同上 \\
6.5.6 & 调制方式 & QPSK\textasciitilde 4096QAM Gray 映射 & \texttt{6.5.6.tex} \\
6.5.7 & 比特加扰 & $\tilde{t}_i{=}t_i\oplus c(i)$ & \texttt{6.5.7-6.5.9.tex} \\
6.5.8 & SC-FDMA & DFT 预编码 $1/\sqrt{M}$ & 同上 \\
6.5.9 & ZC/组跳 & DMRS 基序列 + $u$ 跳变 & 同上 \\
\bottomrule
\end{longtable}

\section{与相关章节的关系}

\begin{longtable}{L{2.2cm}L{1.8cm}L{9.5cm}}
\caption{与相关章节的关系\label{tab:65-rel}} \\
\toprule
\textbf{相关节号} & \textbf{关系} & \textbf{说明} \\
\midrule
6.1 & 上游 & 比特序、链路命名、物理层服务 \\
6.2.2 & 上游 & 帧/符号/$N_{\mathrm{CP}}$/RE 网格 \\
6.2.3\textasciitilde 6.2.6 & 上下游 & 信号/控制/数据/编码消费或供给 6.5 \\
6.3.3 & 下游 & 深覆盖：加扰$\rightarrow$调制$\rightarrow$SC-FDMA；DMRS$\rightarrow$ZC \\
6.6 & 下游 & PMS CSI/时间差/AoA \\
第 8 章 & 下游 & $P_{\mathrm{MAX}}$、射频指标 \\
\bottomrule
\end{longtable}

\section{整章端到端流水线}

\begin{center}
\small
\textbf{控制面}：6.2.4 组装 $\rightarrow$ \textbf{6.5.1} CRC $\rightarrow$ 6.2.6 极化 $\rightarrow$ \textbf{6.5.6} QPSK $\rightarrow$ 映射\\[0.4em]
\textbf{参考信号}：$(n,l,N_{\mathrm{ID}},N_{\mathrm{CP}})$ $\rightarrow$ \textbf{6.5.2}$\rightarrow$\textbf{6.5.3}；或 \textbf{6.5.9.1}$\rightarrow$\textbf{6.5.9} $\rightarrow$ 映射\\[0.4em]
\textbf{数据面}：TB $\rightarrow$ \textbf{6.5.1} $\rightarrow$ 编码 $\rightarrow$ \textbf{6.5.7}（6.5.2）$\rightarrow$ \textbf{6.5.6} $\rightarrow$ \textbf{6.5.4}（←6.5.5.1）$\rightarrow$ 映射\\[0.4em]
\textbf{深覆盖}：\textbf{6.5.7}$\rightarrow$\textbf{6.5.6}$\rightarrow$\textbf{6.5.8}$\rightarrow$映射；DMRS∥\textbf{6.5.9}\\[0.4em]
\textbf{测量}：训练 $\rightarrow$ \textbf{6.5.5.1/2} $\rightarrow$ RSRQ/SINR；CBRA；PMS$\rightarrow$6.6
\end{center}

\textbf{共享底座}：6.5.2 Gold 同时服务 6.5.3、6.5.7、6.5.9.1——同一库、三种 $c_{\mathrm{init}}$，禁止混流。

\newpage
%======================================================================
\part{6.5 物理层通用功能（工程解读）}
%======================================================================

\section{协议章节对应}

本节对应 T/XS 10001-2025 第 6 章第 6.5 节全部子节 6.5.1\textasciitilde 6.5.9（含 6.5.5.x、6.5.6.x、6.5.9.1）。

\section{功能定义}

\begin{enumerate}[nosep]
  \item \textbf{提供} 跨链路复用的 CRC、伪随机序列、参考 QPSK；
  \item \textbf{提供} T 链路开环功控与通信/位置测量量；
  \item \textbf{提供} 比特到星座的调制映射；
  \item \textbf{提供} 加扰、SC-FDMA 预编码与 ZC/组跳；
  \item \textbf{约定} 各块以显式条件调用，不引入章级状态机。
\end{enumerate}

\section{模块边界（上下游及职责划分）}

\subsection{上游模块}

\begin{longtable}{L{4cm}L{5cm}L{5cm}}
\caption{章级上游输入\label{tab:65-up}} \\
\toprule
\textbf{上游} & \textbf{输入} & \textbf{进入 6.5} \\
\midrule
6.2.4/6.2.5 组装 & 信息比特、业务类型 & CRC/加扰/调制阶 \\
XRC/系统消息 & $P_o$、参考功率、ID & 功控/组跳/$c_{\mathrm{init}}$ \\
调度 & $M,N,K$、MCS、$m,\alpha$ & 功控/调制/ZC \\
射频前端 & 频域样本、阵列快照 & 测量 \\
\bottomrule
\end{longtable}

\subsection{下游模块}

\begin{longtable}{L{4cm}L{5cm}L{5cm}}
\caption{章级下游输出\label{tab:65-down}} \\
\toprule
\textbf{下游} & \textbf{输出} & \textbf{说明} \\
\midrule
6.2.3/6.3.3 映射 & 符号/序列/功率缩放 & RE 填充 \\
6.2.6 编码 & CRC 后比特 & 控制/码块 \\
MAC/RRC & 测量上报 & 自适应/LBT \\
6.6 & PMS 量 & 测距定位 \\
发射机 & $\hat{P}_k$ & T 符号发射 \\
\bottomrule
\end{longtable}

\subsection{本模块不负责的内容}

\begin{longtable}{L{5.5cm}L{8.5cm}}
\caption{章级不负责内容\label{tab:65-out}} \\
\toprule
\textbf{不负责} & \textbf{负责节} \\
\midrule
PIB 字段语义与调度策略 & 6.2.4/6.2.5/MAC \\
IFFT/CP/上变频完整射频链 & 6.2.9/6.2.10/6.3.3.6 \\
RTT 最终距离解算 & 6.6 \\
闭环 TPC 累积 & 非 6.5.4 默认范围 \\
\bottomrule
\end{longtable}

\section{在 SLB 通信流程中的角色}

6.5 是 120\,kHz PHY 的\textbf{算法底座}：6.2/6.3 各链路只传参调用，不复制 CRC/Gold/调制/功控实现。接收侧对称使用 CRC 校验、解扰、本地再生参考与测量量。

%----------------------------------------------------------------------
\section{关键参数与强制要求（章级提炼）}
%----------------------------------------------------------------------

\begin{longtable}{L{2.4cm}L{11cm}}
\caption{各子节关键公式/规则速查\label{tab:65-key}} \\
\toprule
\textbf{节} & \textbf{关键点} \\
\midrule
6.5.1 & 五种 $g(D)$；系统码；$L{=}24$ 可选掩码 \\
6.5.2 & $c(i)=(x_1(i{+}1600)+x_2(i{+}1600))\bmod 2$；可预计算 $x_1$ \\
6.5.3 & $c_{\mathrm{init}}=2^{10}(15(n{+}1)+l{+}1)(4N_{\mathrm{ID}}{+}1)+4N_{\mathrm{ID}}+N_{\mathrm{CP}}$；$M_{\mathrm{PN}}{=}322$ \\
6.5.4 & $P=10\log_{10}M+P_o+P_L$；$P_L{=}$参考功率$-$RSRP；超 $P_{\mathrm{MAX}}$ 均分 \\
6.5.5 & RSRQ/SINR 同符号；CBRA 门限 $-87/-62\,\mathrm{dBm}$ \\
6.5.6 & 公式整体负号；数据须先加扰；与 6.5.3 符号约定不同 \\
6.5.7 & 默认 $c_{\mathrm{init}}=n_f\cdot 2^{24}+N_{\mathrm{G-PID}}$ \\
6.5.8 & $M_{\mathrm{symb}}$ 整除 $M$；强制 $1/\sqrt{M}$ \\
6.5.9 & $m{=}2$ 查表；组跳 $u=(f_{\mathrm{gh}}+f_{\mathrm{SS}})\bmod 30$ \\
\bottomrule
\end{longtable}

\begin{longtable}{L{3.5cm}L{2.5cm}L{8cm}}
\caption{三类 Gold $c_{\mathrm{init}}$（禁止混用）\label{tab:65-gold3}} \\
\toprule
\textbf{场景} & \textbf{节} & \textbf{初始化} \\
\midrule
参考 QPSK & 6.5.3 & $n,l,N_{\mathrm{ID}},N_{\mathrm{CP}}$ \\
比特加扰 & 6.5.7 & 业务表（数据/广播/510/高层） \\
ZC 组跳 & 6.5.9.1 & $\lfloor n_{\mathrm{ID}}^X/30\rfloor$ \\
\bottomrule
\end{longtable}

%----------------------------------------------------------------------
\section{本节核心详解}
%----------------------------------------------------------------------

\subsection{算法库切分建议}

\begin{longtable}{L{4.5cm}L{9.5cm}}
\caption{工程库切分\label{tab:65-libs}} \\
\toprule
\textbf{库} & \textbf{覆盖} \\
\midrule
\texttt{slb\_phy\_common} & 6.5.1\textasciitilde 6.5.3 \\
\texttt{slb\_phy\_pwr\_meas} & 6.5.4\textasciitilde 6.5.5 \\
\texttt{slb\_phy\_mod} & 6.5.6（含 LUT） \\
\texttt{slb\_phy\_scram/scfdma/zc} & 6.5.7\textasciitilde 6.5.9（Gold 回调 common） \\
\bottomrule
\end{longtable}

\subsection{关键耦合}

\begin{itemize}[nosep]
  \item 功控依赖 RSRP（6.5.5.1 $\rightarrow$ 6.5.4）；
  \item 数据调制依赖加扰（6.5.7 $\rightarrow$ 6.5.6）；
  \item 深覆盖在调制后插入 SC-FDMA（6.5.6 $\rightarrow$ 6.5.8）；
  \item 组跳与加扰均调 Gold，但 $c_{\mathrm{init}}$/流独立。
\end{itemize}

%----------------------------------------------------------------------
\section{数据类型、长度与维度}
%----------------------------------------------------------------------

\begin{longtable}{L{3.5cm}L{4cm}L{6.5cm}}
\caption{章级数据对象摘要\label{tab:65-dim}} \\
\toprule
\textbf{对象} & \textbf{典型维度} & \textbf{备注} \\
\midrule
CRC 比特流 & $A$ / $A{+}L$ & $L\in\{6,8,16,24\}$ \\
Gold / 加扰 & $M_{\mathrm{PN}}$、$N_{\mathrm{bit}}$ & 调用方定长 \\
RS QPSK & 161 & DC $k{=}80$ \\
调制符号 & $N_{\mathrm{bit}}/m$ & $m\in\{2,4,6,8,10,12\}$ \\
功控 & 标量 / $K$ 子载波 & 限幅后 $\hat{P}_k$ \\
ZC & $M_{\mathrm{SC}}^{\mathrm{RS}}$ & $u{=}0\ldots 29$ \\
\bottomrule
\end{longtable}

%----------------------------------------------------------------------
\section{时序关系与触发条件}
%----------------------------------------------------------------------

\begin{longtable}{L{3.2cm}L{5.8cm}L{5cm}}
\caption{章级触发条件\label{tab:65-trig}} \\
\toprule
\textbf{链路} & \textbf{进入条件} & \textbf{失败回退} \\
\midrule
控制发送 & 组装完成 & CRC/极化参数错则停 \\
数据发送 & MCS/$c_{\mathrm{init}}$/$M$ 已知 & 加扰或调制校验失败则停 \\
T 发射功控 & RSRP 与 $P_o$ 就绪 & $M{=}0$ 或缓冲不足报错 \\
深覆盖 & $M_{\mathrm{symb}}$ 可整除 & 不整除禁止 DFT \\
测量上报 & 同符号样本齐 & 跨符号算 RSRQ/SINR 拒绝 \\
\bottomrule
\end{longtable}

%----------------------------------------------------------------------
\section{处理步骤}
%----------------------------------------------------------------------

\subsection{端到端 A：控制块}

\textbf{Step 1}：组装 $a[]$ $\rightarrow$ 选 CRC 多项式。\\
\textbf{Step 2}：6.5.1 编码 $\rightarrow$ 6.2.6 极化。\\
\textbf{Step 3}：6.5.6.1 QPSK $\rightarrow$ 映射。\\
失败任一步立即返回，不进入``下一状态''。

\subsection{端到端 B：数据 TB}

\textbf{Step 1}：TB CRC（6.5.1）$\rightarrow$ 编码。\\
\textbf{Step 2}：6.5.7 加扰（6.5.2）。\\
\textbf{Step 3}：6.5.6 调制。\\
\textbf{Step 4}：若 6.3.3 则 6.5.8；否则直接映射。\\
\textbf{Step 5}：6.5.5.1$\rightarrow$6.5.4 功率后发射。

\subsection{端到端 C：参考 / DMRS}

\textbf{条件}：若伪随机 RS $\rightarrow$ 6.5.2+6.5.3；若 ZC DMRS $\rightarrow$ 6.5.9.1（可选）+6.5.9 $\rightarrow$ OCC/$β$。\\
\textbf{失败}：$m{=}2$ 误用闭式或 $c_{\mathrm{init}}$ 混用则判不合格。

\subsection{端到端 D：测量}

同符号 RSRP+RSSI $\rightarrow$ RSRQ/SINR；独立触发 CBRA/PMS。无测量态机，仅检查样本是否齐备。

%----------------------------------------------------------------------
\section{协议中允许改变的参数（章级）}
%----------------------------------------------------------------------

\begin{longtable}{L{4.5cm}L{4cm}L{5.5cm}}
\caption{章级可配置参数\label{tab:65-tunable}} \\
\toprule
\textbf{参数族} & \textbf{来源} & \textbf{影响面} \\
\midrule
CRC 类型 / 掩码 & 信息块类型 & 6.5.1 \\
$c_{\mathrm{init}}$ / $M_{\mathrm{PN}}$ & 调度/ID & 6.5.2/3/7/9.1 \\
$P_o$ / $P_{\mathrm{MAX}}$ / $M$ & XRC/射频/调度 & 6.5.4 \\
MCS / 调制阶 & 6.2.5 & 6.5.6 \\
$M_{\mathrm{SC}}^{\mathrm{Total}}$ / $m,\alpha$ & 带宽/DMRS & 6.5.8/9 \\
\bottomrule
\end{longtable}

%----------------------------------------------------------------------
\section{链路调用总表}
%----------------------------------------------------------------------

\begin{longtable}{L{3.8cm}L{2.2cm}cccc}
\caption{典型链路与 6.5 子节调用\label{tab:65-link}} \\
\toprule
\textbf{链路} & \textbf{位置} & \textbf{1--3} & \textbf{4--5} & \textbf{6} & \textbf{7--9} \\
\midrule
SAB/控制 & 6.2.4 & CRC & --- & QPSK & --- \\
第一/二类数据 & 6.2.5 & CRC & 功控/SINR & MCS 阶 & 加扰 \\
SRS/DMRS(QPSK) & 6.2.3 & Gold+QPSK & --- & --- & --- \\
T 接入/ACK & 6.2.4 & --- & 功控 & --- & --- \\
深覆盖广播/控制 & 6.3.3 & --- & --- & 调制 & 加扰+SC-FDMA \\
深覆盖 DMRS & 6.3.3.2.2 & Gold(组跳) & --- & --- & ZC \\
选网/LBT/定位 & 6.5.5/6.6 & --- & 测量 & --- & --- \\
\bottomrule
\end{longtable}

%----------------------------------------------------------------------
\section{约束与实现要点}
%----------------------------------------------------------------------

\begin{enumerate}[nosep]
  \item \textbf{【分册为准】}：公式/接口细节以四份分册为准，本整合文档不另立冲突条款；
  \item \textbf{【Gold 三场景隔离】}：6.5.3 / 6.5.7 / 6.5.9.1 的 $c_{\mathrm{init}}$ 与流偏移禁止混用；
  \item \textbf{【1600 偏移】}：Gold 输出必须用 $i{+}1600$，不可省略；
  \item \textbf{【数据先加扰后调制】}：6.5.6 数据输入必须来自 6.5.7；
  \item \textbf{【调制负号 / RS 正号】}：6.5.6.1 与 6.5.3 QPSK 符号约定不同，禁混 API；
  \item \textbf{【同符号测量】}：RSRQ/SINR 的 RSRP/RSSI 必须同符号；
  \item \textbf{【线性域平均】}：功率先线性平均再转 dB；
  \item \textbf{【开环功控】}：6.5.4 不做 TPC 累积；
  \item \textbf{【$m{=}2$ 查表】}：ZC 短序列禁止误用闭式；
  \item \textbf{【可重入】}：各库无静态全局可变状态，缓冲由调用方提供；
  \item \textbf{【无状态机】}：整章实现用事件/配置/标志/计数器等显式条件推进（样本是否就绪、$N_{\mathrm{bit}}\bmod m$、是否深覆盖、$N{>}1$ 是否限幅、是否组跳等）；禁止维护章级或库级 IDLE/WORK/DONE FSM，失败立即返回调用方。
\end{enumerate}

%----------------------------------------------------------------------
\section{与标准其他章节的衔接}
%----------------------------------------------------------------------

\begin{itemize}[nosep]
  \item \textbf{6.2.x}：帧结构、信号/控制/数据、编码是 6.5 的主要调用方；
  \item \textbf{6.3.3}：深覆盖强制串联 6.5.7/6.5.6/6.5.8 与 6.5.9；
  \item \textbf{6.6}：位置流程消费 6.5.5.6\textasciitilde 6.5.5.8；
  \item \textbf{分册}：\texttt{6.5.1-6.5.3.tex}、\texttt{6.5.4-6.5.5.tex}、\texttt{6.5.6.tex}、\texttt{6.5.7-6.5.9.tex}。
\end{itemize}

\vfill
\begin{center}
\small\textcolor{gray}{精确参数与字段定义以 T/XS 10001-2025 正式标准及各分册为准；\\
本文档提供 6.5 章整合工程入口与端到端约束，供 SLB 实现与联调参考。}
\end{center}

\end{document}
